\section{Record-only analysis without supplied call endpoints}\label{supp:no-call}
This extension reuses all 3,000 main episodes. The analyst fixed its code and candidate grid before computing new test scores, with the older task outcomes already known. It is retrospective reused-workload evidence, not a new cohort. The underlying live collection, primary choices and tests remain unchanged.

\subsection{Inputs and observation scope}
Each episode contributes seven final, closed connection logs, linked by the episode's connection bundle. All outer-type23 records are included; record time is made relative to the earliest such record. Initialization and closing records are retained. The feature whitelist is relative time, destination, direction and ciphertext record length. Labels and split identifiers join evaluation rows but never enter features. TEST inference reads no per-call terminal, operation, private role or HTTP stream mapping. The later supervised extension uses TRAIN/validation role annotations only for fitting and selection (Section~\ref{supp:record-context}). A source-hash manifest binds every input. Whole-task association is retained, so this is removal of call endpoints, not task localization.

\subsection{Candidate representations and fitting}
Within each destination connection, records separated by more than 5, 20, 50, 100 or 250\,ms form separate proposed events. Events retain the first and last timestamps and directional ciphertext sums. Destination summaries include event count, directional bytes, record count and maximum event volume; timing variants also include first/last times and total event duration. Ordered features use the first four and last sixteen events, with destination one-hot and log byte counts, adding duration and preceding gap in timing variants. Raw-record candidates treat each record as an event and use the same summaries. No plaintext or true call interval is required.

LR uses standardization, $C=1$ and at most 3,000 iterations. RF has 300 trees. The fixed seed is 111121 and numerical libraries use two threads. Per-condition training fits each representation. A two-component mixture fits log gaps from that condition's training events without role labels. The crossing between its short/long components defines removal; the first event is retained as a heuristic. The remaining events go to a no-cover-trained task LR/RF at the same segmentation/view. Validation chooses baseline, removal and overall attacks by task accuracy with fixed-order ties. The selection file is written before that cell's test predictions.

There are forty-four fixed-grid candidates per model/condition: five gaps, two views, two learners for refitted baseline and clean-model removal, plus four raw-record candidates. Each cell retains 120 training, 60 validation and 120 test fixtures. The primary comparisons are selected removal minus selected baseline on immediate cover, with exact paired McNemar and Holm across the two models. Ten thousand paired test-fixture bootstrap draws give conditional intervals; they omit model selection, refitting and workload-selection uncertainty. Other arms and test maxima are descriptive. A test maximum is neither an independently selected confirmatory attack nor an upper bound on all possible inference.

\subsection{Exploratory protocol-aware eventization}
After the Qwen idle-grid test results were known, a separate plan fixed one public-protocol alternative. An upward encrypted record of at least 4,096 bytes starts a proposed transaction. The first two such requests per connection correspond to the public MCP initialize/initialized setup sequence and are removed from this representation. A transaction ends at the last downward record of at least 4,096 bytes before the next trigger, or at its start if no such response exists. This rule exploits the configured padding quantum and same-service serialization; it reads no true operation, call endpoint or private role. It is not an HTTP/2 plaintext parser or a general multiplexing solution.

The same size/timing views, LR/RF and clean-model GMM removal yield eight further candidates per cell. TRAIN/validation/test remain fixed, and validation selections are saved before their test predictions. These are exploratory results with no new superiority family. The first fixed-grid results and tests are preserved in full. The displayed Max spans all fifty-two candidates, while B/R still refer to the original fixed-grid validation choices. Complete candidate counts and predictions are in the local analysis package.

\subsection{Exploratory context removal and joint decoding}\label{supp:record-context}
This extension was designed after the preceding record-only results were known. It keeps all 3,000 original episode observations byte-identical, retains the fixed 120/60/120 splits, and adds learned context attacks on the same five idle gaps and protocol eventizer. Each cell reports 76 candidates: 28 refitted baselines (six eventizations, two views, two learners, plus four raw-record classifiers), 24 context-removal candidates and 24 joint-decoding candidates. The older GMM removals remain separately reported in the original 52-candidate results below; they are not silently replaced or counted twice.

An event's 32 detector features describe duration, preceding and same-service gaps, destination and predecessor, position and first/last flags, directional bytes, record count, elapsed time, next-event context and destination recurrence. TRAIN and validation calls label an event as cover when it overlaps any controller-cover or shadow call on the same destination, allowing 2\,ms slack. Such an event may also contain honest data. LR/RF detectors use balanced class weights; validation role F1, precision and fixed-order ties select the detector and a threshold in .05 increments from .05 to .95. These role labels are not TEST inference inputs.

Context removal drops flagged events and applies a no-cover-trained task classifier. The first event is retained as a heuristic. Joint decoding considers removal of zero, one or two events among up to ten highest-probability candidates, excluding the first event. Its fixed removal-count weights are .02, .49 and .49. A MAP rule uses the detector score plus task confidence, with weight in $\{0,.5,1,2,4\}$; marginal decoding is also a candidate. Validation task accuracy selects that internal rule and the final candidate, with fixed-order ties. Each task representation uses the original standardized LR or 300-tree RF, seed 111121 and two threads. A maximum of two inferred events is not equivalent to a two-call budget.

The selected context/joint family is compared with the selected baseline in each of eight cover cells, using two-sided exact paired McNemar testing and one eight-row Holm family. Table~\ref{tab:supp-record-tests} gives every discordance and test. Five comparisons improve significantly: all three Qwen controller cells and both Llama history cells. Llama immediate has an uncorrected $p=.0501$ and corrected $p=.1503$; both shadow contrasts are negative and nonsignificant. Overall validation selection across baseline and added attacks keeps the 69/120 baseline for each shadow. All 76 candidate validation/test counts, selected names and descriptive maxima are in the aggregate artifact; test maxima do not replace validation choices.

All validation and test task predictions are saved and sealed before loading TEST private call truth for the event audit. Table~\ref{tab:supp-record-audit} reports the eventization selected by each strong attack; the released audit covers all six eventizers. Cover-label P/R/F1, call-boundary P/R/F1 at IoU .5, and mixed/pure/unmatched event counts remain separate. Mixed positive events can contain honest data, so a role F1 does not certify correct call removal. Neither audit endpoint is a test-time feature or a selection criterion. Whole-episode connection association is still supplied; this experiment does not discover tasks in an unassociated stream.
\begin{table*}[t]\centering\footnotesize
\caption{All eight exploratory context/joint minus baseline tests on record-only observations. Discordances and denominators are paired fixture counts; $n=120$ in each cell. Holm correction covers all eight rows.}\label{tab:supp-record-tests}
\setlength{\tabcolsep}{3pt}
\begin{tabularx}{\textwidth}{@{}llrrrrr@{}}\toprule
Planner & Policy & J-only & B-only & Difference (pp) & Raw $p$ & Holm $p$ \\\midrule
Qwen & Immediate & 33 & 1 & +26.7 & $4.1\!\times\!10^{-9}$ & $2.9\!\times\!10^{-8}$ \\
Qwen & History, seq. & 25 & 8 & +14.2 & 0.0046 & 0.0182 \\
Qwen & History, conc. & 39 & 15 & +20.0 & 0.0015 & 0.0079 \\
Qwen & Shadow & 20 & 29 & -7.5 & 0.2529 & 0.5057 \\
Llama & Immediate & 22 & 10 & +10.0 & 0.0501 & 0.1503 \\
Llama & History, seq. & 26 & 7 & +15.8 & 0.0013 & 0.0079 \\
Llama & History, conc. & 35 & 1 & +28.3 & $1.1\!\times\!10^{-9}$ & $8.6\!\times\!10^{-9}$ \\
Llama & Shadow & 14 & 19 & -4.2 & 0.4869 & 0.5057 \\
\bottomrule\end{tabularx}\end{table*}

\begin{table*}[t]\centering\footnotesize
\caption{Test audits for the eventization used by the validation-selected strong attack. Role P/R/F1 labels an event as cover if it overlaps any cover call; boundary P/R/F1 matches events to calls at IoU .5. Mixed events overlap real and cover calls. All test predictions were saved before these private-truth audits. The complete six-eventizer audit is in the aggregate artifact.}\label{tab:supp-record-audit}
\setlength{\tabcolsep}{3pt}
\begin{tabularx}{\textwidth}{@{}llrYYrr@{}}\toprule
Planner & Policy & Event gap & Role P/R/F1 & Boundary P/R/F1 & Mixed & Events \\\midrule
Qwen & Immediate & 100 ms & 0.994/0.989/0.992 & 0.411/0.938/0.571 & 14 & 1534 \\
Qwen & History, seq. & 100 ms & 0.784/0.870/0.825 & 0.420/0.969/0.586 & 0 & 1547 \\
Qwen & History, conc. & 100 ms & 0.807/0.934/0.866 & 0.366/0.896/0.520 & 7 & 1643 \\
Qwen & Shadow & Protocol & 0.505/0.912/0.650 & 0.390/0.612/0.476 & 0 & 1533 \\
Llama & Immediate & 100 ms & 1.000/1.000/1.000 & 0.475/0.947/0.632 & 17 & 1629 \\
Llama & History, seq. & 250 ms & 0.963/0.857/0.907 & 0.491/0.999/0.659 & 0 & 1659 \\
Llama & History, conc. & 250 ms & 0.988/0.929/0.958 & 0.488/0.982/0.652 & 10 & 1649 \\
Llama & Shadow & 5 ms & 0.488/0.998/0.655 & 0.000/0.000/0.000 & 0 & 8629 \\
\bottomrule\end{tabularx}\end{table*}


\subsection{Complete original results}
\begin{table*}[t]\centering\small
\caption{Selected record-only attacks. B/R/A refer to validation-selected fixed-grid baseline, removal and overall attackers. Proto B/R gives the separate exploratory protocol-aware selection; Max spans both suites. Counts are out of 120. Complete candidate outcomes remain in the local analysis package.}\label{tab:no-call-results}
\begin{tabularx}{\textwidth}{@{}llrrrrY@{}}\toprule Planner & Policy & B & R & Proto B/R & Max & Selected fixed-grid attacks (B; R; A) \\\midrule
Qwen & No cover & 115 & 109 & 117/93 & 117 & events 5ms LR; remove 20ms timing RF; events 5ms LR \\
Qwen & Immediate & 80 & 69 & 68/91 & 91 & records timing RF; remove 5ms RF; records timing RF \\
Qwen & History, sequential & 73 & 74 & 66/45 & 75 & records timing RF; remove 5ms RF; remove 5ms RF \\
Qwen & History, concurrent & 56 & 61 & 51/47 & 62 & events 5ms LR; remove 5ms RF; remove 5ms RF \\
Qwen & Planner shadow & 69 & 65 & 66/52 & 73 & events 20ms RF; remove 100ms timing RF; events 20ms RF \\
Llama & No cover & 108 & 101 & 109/88 & 112 & events 5ms timing LR; remove 20ms RF; events 5ms timing LR \\
Llama & Immediate & 90 & 86 & 90/87 & 97 & records RF; remove 5ms timing RF; records RF \\
Llama & History, sequential & 88 & 56 & 90/71 & 93 & events 5ms timing RF; remove 20ms LR; events 5ms timing RF \\
Llama & History, concurrent & 78 & 62 & 83/68 & 88 & events 50ms RF; remove 5ms LR; events 50ms RF \\
Llama & Planner shadow & 69 & 56 & 68/52 & 69 & events 5ms RF; remove 5ms RF; events 5ms RF \\
\bottomrule\end{tabularx}\end{table*}

\begin{table*}[t]\centering\small
\caption{Prespecified record-only immediate-cover comparisons. Discordance counts are removal-only/baseline-only correct. Intervals condition on fitted predictors and fixtures.}\label{tab:no-call-tests}
\begin{tabularx}{\textwidth}{@{}lrrrrYrr@{}}\toprule Planner & Removal & Baseline & $+/-$ & Difference pp & Conditional 95\% interval & Raw $p$ & Holm $p$ \\\midrule
Qwen & 69 & 80 & 3/14 & -9.2 & [-15.8, -2.5] & 0.01273 & 0.02545 \\
Llama & 86 & 90 & 10/14 & -3.3 & [-11.7, 4.2] & 0.5413 & 0.5413 \\
\bottomrule\end{tabularx}\end{table*}

\begin{table*}[t]\centering\small
\caption{Independent call-boundary audit, IoU .5, after prediction. Each row uses120 associated test episodes. Idle selects its burst gap by validation task accuracy; Protocol uses public padding/setup rules. Values are precision / recall / F1, including spurious proposed events. These truth comparisons are not attacker inputs.}\label{tab:no-call-boundaries}
\begin{tabularx}{\textwidth}{@{}llYY@{}}\toprule Planner & Policy & Idle-gap eventization & Protocol-aware eventization \\\midrule
Qwen & No cover & 0.000/0.000/0.000 & 0.233/0.434/0.303 \\
Qwen & Immediate & 0.000/0.000/0.000 & 0.229/0.412/0.295 \\
Qwen & History sequential & 0.000/0.000/0.000 & 0.272/0.489/0.350 \\
Qwen & History concurrent & 0.000/0.001/0.000 & 0.266/0.478/0.342 \\
Qwen & Shadow & 0.012/0.078/0.021 & 0.390/0.612/0.476 \\
Llama & No cover & 0.000/0.000/0.000 & 0.294/0.510/0.373 \\
Llama & Immediate & 0.000/0.000/0.000 & 0.302/0.512/0.380 \\
Llama & History sequential & 0.000/0.000/0.000 & 0.378/0.641/0.476 \\
Llama & History concurrent & 0.357/0.802/0.494 & 0.358/0.606/0.450 \\
Llama & Shadow & 0.000/0.000/0.000 & 0.428/0.637/0.512 \\
\bottomrule\end{tabularx}\end{table*}

The original independent boundary diagnostic is computed only after attack predictions, using saved private call truth. It is not an attacker feature or a segmentation selection criterion. The local package retains every held-out prediction and input hash. The public aggregate release supplies candidate validation/test counts, selections and reconstruction code; individual record metadata and predictions remain local.
